% Overleaf figure wrappers for Whisper--fairseq comparison plots.
% Upload/copy the corresponding PDFs from plots_lm/ into your Overleaf project.
% Requires: \usepackage{graphicx}

\begin{figure}[t]
\centering
\includegraphics[width=\linewidth]{whisper_fairseq_two_panel.pdf}
\caption{%
Comparison of Whisper and fairseq degradation under structured label noise.
Panel A reports word error rate (WER) for each noise condition, with color denoting the condition and hatch/marker style denoting model family.
Panel B shows baseline-relative failure behavior, plotting the change in WER against the change in normalized fluency relative to each model family's own clean baseline.
Rightward movement indicates greater lexical degradation, while downward movement indicates reduced language-model fluency.
Whisper results use the 64\% noise setting; fairseq results use the 8\% noise setting with the same Qwen3-0.6B fluency scorer.
}
\label{fig:whisper-fairseq-fluency}
\end{figure}

\begin{figure}[t]
\centering
\includegraphics[width=\linewidth]{whisper_fairseq_two_panel_repetition.pdf}
\caption{%
Comparison of recognition degradation and repetition behavior for Whisper and fairseq.
Panel A repeats the WER comparison across noise conditions.
Panel B replaces fluency with the change in mean trigram repetition count relative to each model family's own clean baseline.
Rightward movement indicates increased lexical error, while upward movement indicates stronger repetitive output behavior.
This separates fluent hallucination-like degradation from repetition-driven failure modes.
}
\label{fig:whisper-fairseq-repetition}
\end{figure}

\begin{figure}[t]
\centering
\includegraphics[width=\linewidth]{whisper_fairseq_repetition_hallucination_patterns.pdf}
\caption{%
Failure-pattern comparison between Whisper and fairseq under structured label noise.
Panel A shows the change in WER against the change in four-gram repetition count, both measured relative to each model family's own clean baseline.
The shaded upper region marks high-repetition behavior, separating repeated-output failures from other recognition errors.
Panel B plots absolute WER against absolute fluency on a 0--1 scale.
The orange band marks fluent plausible sentences, showing that most noisy-condition outputs remain language-model plausible even when WER increases.
The softly cross-hatched high-WER region marks transcription failure, where errors are dominated by lexical mismatch rather than loss of sentence fluency.
Whisper results use the 64\% noise setting; fairseq results use the 8\% noise setting, with fluency averaged across GPT-2 and Qwen3-0.6B scores.
}
\label{fig:whisper-fairseq-failure-patterns}
\end{figure}
